\section{The Bechinah Mathematical Model}
\label{sec:bechina_model}

We define the fundamental object of study, a \emph{Bechinah}, as a function over a three-dimensional domain space.

\begin{definition}[Value, Relation, Scale Spaces]
Let $\mathcal{V}$ denote a value space, $\mathcal{R}$ denote a space of relational operators, and $\mathcal{S} = \mathbb{R}^+$ denote a positive scalar scale space.
\end{definition}

\begin{definition}[Bechinah Function]
A \emph{Bechinah} $B$ is a configuration map:
\begin{equation}
\boxed{B = \operatorname{Bechina}(V, R, S)}
\end{equation}
where $V \in \mathcal{V}$, $R \in \mathcal{R}$, and $S \in \mathcal{S}$.
\end{definition}

In simple linear approximations, one might consider a multiplicative heuristic for informational content:
\begin{equation}
M = V \cdot R \cdot S
\end{equation}
However, such a simple product fails when relations or scale exhibit non-linear bounds or phase transitions. Therefore, we model general meaningful information ($M$) as a general non-linear operator $F$:
\begin{equation}
M = F(V, R, S)
\end{equation}
This formulation preserves mathematical safety against division by zero and accounts for non-linear scale coupling.
